\section{Related work}
\label{sec:related-work}

\subsection{Directed feedback arc set: formulations and algorithmic foundations}

The feedback arc set problem and its weighted variants sit at the intersection of graph optimization, ordering, and ranking. In directed graphs, the problem can be phrased either as deleting a minimum-weight set of arcs to obtain an acyclic graph or as finding an ordering with minimum backward weight. That equivalence makes the problem relevant both to graph-theoretic cycle breaking and to applications that begin from pairwise precedence or preference data. Classical digraph formulations identified minimum feedback arc sets as a central cyclic-structure question \cite{LC66}, early weighted treatments framed the problem as an ordering objective \cite{F90}, and the general decision problem is computationally hard \cite{K72}. The approximation literature is broader for vertex-deletion or undirected variants than for the weighted directed setting considered here, but it still provides important methodological context. \cite{ENSS98} study approximation algorithms for minimum feedback sets and multicuts in directed graphs, while \cite{HGH21} analyze tight localizations of feedback sets in more recent algorithm-engineering work. In particular, \cite{BYGR98} place local-ratio reasoning in a general approximation framework, showing how weighted combinatorial objectives can be decomposed while preserving feasible improvement structure. Their setting is not the same as minimum weighted feedback arc set on directed graphs, and we do not rely on their theorem to obtain a new approximation guarantee here. However, the paper is directly relevant to the design logic of iterative weight reductions on conflict structures.

For directed feedback problems, \cite{DF03} provide the closest methodological antecedent to the cycle-reduction component used here. Their contribution is important for two reasons. First, it treats directed feedback problems algorithmically rather than only as ordering formulations. Second, it shows how local-ratio reductions can be applied to directed cycles in a way that is compatible with practical heuristics. The local-ratio step in the present work is inherited in spirit from that line of research: when a directed cycle is found, all arcs on the cycle are reduced by the minimum weight on the cycle, at least one arc becomes tight, and the active graph is simplified. What is new here is not the local-ratio principle itself, but an engineered sparse-graph implementation that records removed arcs explicitly, reconstructs a deterministic topological order once the active graph becomes acyclic, and then runs a heavy-first add-back step to reduce unnecessary deletions.

A systematic comparison of sorting heuristics for feedback arc set appears in \cite{BH13}, who evaluate Eades-family net-score rules, SCC-based decomposition approaches, and other greedy ordering strategies on directed graph instances. Their study provides useful empirical context for the heuristic family: decomposing the problem along SCC boundaries generally improves ordering quality compared with applying net-score rules to the graph as a whole. The present work uses that structural insight as a design principle in both the seed construction and the refinement stage.

The contribution is therefore best understood as an algorithm-engineering framework built around a known reduction principle. The emphasis is on deterministic data handling, compatibility with sparse weighted digraphs, and integration with a second-stage refinement process that improves only when the global objective decreases.

\subsection{Heuristics, library implementations, and order-local search}

Heuristic work on feedback arc set and related ordering tasks spans greedy graph rules, local search, and software packages that expose ordering or arc-deletion primitives. A classical baseline is the greedy heuristic of \cite{ELS93}. Their method is attractive because it is fast, interpretable, and structurally simple: it repeatedly extracts sources, sinks, or high net-score vertices to construct an order. An earlier denser-graph variant in the same heuristic family appears in \cite{EL95}. Even when a modern study ultimately outperforms such a heuristic, it remains an essential reference point because many later practical systems either implement that rule directly or build on the same score-based intuition. In the present paper, Eades-style ideas appear in two places: as the historical baseline that motivates fast greedy ordering, and as the conceptual starting point for external and adapted software baselines used in the experiments.

The external-baseline design intentionally mixes library-grade and specialized open-source tools. The python-igraph documentation for \texttt{Graph.feedback\_arc\_set} \cite{python_igraph_feedback_arc_set} exposes both an Eades-style heuristic and exact integer-programming modes. It is a useful baseline because it reflects what a practitioner could call from a general-purpose graph library. In contrast, the DRMacIver/FAS tool \cite{drmaciver_feedback_arc_set} is a matrix-based pairwise-ordering heuristic rather than a general graph-analysis package. Its documentation formulates the input as a nonnegative \(n \times n\) matrix \(W\), where \(W_{ij}\) represents evidence that item \(i\) should precede item \(j\), and seeks a permutation maximizing the total forward weight \(\sum_{p_i < p_j} W_{ij}\). The repository describes the method as deterministic and reports a local-optimality property with respect to single-element moves; it does not present the tool as an exact optimizer. Because the documented model is matrix-based and has \(O(n^2)\) complexity in the number of items, it is especially natural as a dense pairwise-ordering baseline, while remaining usable in our experiments through the sparse entry-list interface. Empirically, it is the strongest external baseline on the sparse standard benchmark and the dominant method on the dense LOLIB transfer test.

The remaining baselines in the experiments are deliberately simple and interpretable. Weighted net-score ordering and random multistart are not intended to represent state of the art; they provide calibration points. Net-score and win-based tournament rules are classical ranking ideas with theoretical support in the weighted tournament literature \cite{CFR10}. The weighted Eades adaptation used here should also be interpreted carefully. The original Eades--Lin--Smyth paper \cite{ELS93} does not define the weighted adaptation used in this study, so that baseline is described as a local adaptation rather than as the original authors' algorithm. That distinction is part of the methodological hygiene of the study: external baselines should be credited as external when they are external, and local adaptations should be labeled as such even when they are inspired by well-known methods.

Beyond greedy seeds, simpler order-local search strategies provide useful methodological context. Adjacent-swap local search and single-vertex insertion local search represent the class of moves that improve an ordering by reordering one or two elements at a time without reference to SCC structure. A natural question for any SCC-aware refinement framework is whether such moves, seeded from a strong constructive heuristic, can replicate its gains. The study includes a direct comparison in EXP7 to determine whether these order-local moves suffice, allowing the contribution of SCC-specific structure to be isolated from generic greedy improvement.

There is a broader family of learning-based ranking methods that could be discussed without entering the experimental comparison set. GNNRank \cite{GNNRank22}, for example, frames ranking from pairwise comparisons as a trainable graph-neural ranking problem. It is relevant as a signal that the ranking literature has moved beyond purely combinatorial heuristics, but it is not used as a baseline here because the present study focuses on deterministic, training-free graph heuristics rather than learned ranking models. Recent weighted feedback arc set heuristics such as the minimal-and-stable construction of Cavallaro and Cutello \cite{CC25} are likewise outside the comparison set because they were not rerun for this study, but they illustrate continued heuristic interest in weighted feedback arc deletion.

\subsection{Exact methods and computational hardness}

The exact side of the literature serves two different roles in this paper. First, it sets expectations about computational hardness. Second, it provides the reference standard for the small-instance validation experiment. \cite{BSNA21} give a modern exact treatment of minimum feedback arc set through formulations that connect the problem to broader combinatorial optimization machinery. Their work is not the computational baseline for the full sparse benchmark in this paper; the sparse benchmark is too large for exact optimization to be the main comparison mode. Instead, exact methods enter the study through a limited validation regime on small instances, where the question is not whether exact optimization scales globally, but whether the heuristic framework is near-optimal when exact answers are available.

This distinction also helps separate feedback arc set from the larger family of ordering models. Dense linear ordering formulations are often written as forward-weight maximization on complete matrices, whereas sparse directed graph formulations arise more naturally as backward-weight minimization on incomplete digraphs. The two are mathematically linked, but they are not computationally identical in practice.

The present study deliberately moves between the two views. The sparse benchmark is the main claim-bearing environment and is presented in graph terms. The LOLIB transfer test is presented in linear-ordering terms because those instances are dense complete ordering problems. Keeping both formulations visible is important. Otherwise the reader could mistakenly treat good sparse-graph performance as evidence of dense linear-ordering dominance, which the experiments do not support.

At larger scales, exact computation is often replaced by decomposition, approximation, or specialized engineering. \cite{SST16} provide one useful point of reference by treating feedback arc set computation in a web-scale setting. Tournament and dense-ordering settings have their own fast heuristic literature, including local-search methods for weighted tournament feedback arc set \cite{FLRS10} and fixed-parameter approaches for the unweighted tournament case \cite{ALS09}. Their work is not directly comparable to the exact small-instance validation in this paper, but it illustrates the broader computational reality: scalable feedback arc methods are usually driven by structure, heuristics, or deployment constraints rather than by globally exact optimization. The goal here is not to displace exact methods; it is to present a heuristic framework whose computational behavior is strong on the sparse benchmark and whose limitations on dense complete ordering instances are reported explicitly.

\subsection{Dense linear ordering and tournament methods as a structurally distinct regime}

Dense linear ordering and tournament methods constitute a structurally distinct computational regime from the sparse weighted-digraph setting studied here, and the dense-transfer side of the paper is best understood in that context. The dense-transfer side of the paper is best understood through the linear ordering literature. \cite{MRD12} provide a benchmark-oriented account of the Linear Ordering Problem Library (LOLIB) and its heuristic evaluation culture. That paper is the most appropriate formal reference for LOLIB-style benchmark usage here because it anchors the library in a published benchmark-comparison setting rather than only as a download page. The official LOLIB page itself \cite{lolib_library} remains important for provenance and dataset identification, especially because benchmark users often obtain the instances directly from library pages or their mirrors. Together, these sources explain why dense complete tournaments have a different methodological status from sparse DIMACS-style directed graphs: they are not merely larger graphs of the same type, but instances drawn from a different benchmark tradition with different algorithmic expectations.

The sparse benchmark side has a different provenance. The graph-benchmarks collection \cite{graph_benchmarks_repo} is a public set of directed instances in DIMACS-style format rather than a linear-ordering library.

That difference supports a central design choice: the sparse and dense benchmarks test different claims. The sparse benchmark is the main evidence for practical value on nonnegative weighted digraphs. The dense benchmark is a transfer test used to determine where a general weighted-digraph heuristic stops being competitive against dense-ordering-native methods.

The ranking-from-pairwise-comparisons viewpoint connects these benchmark families conceptually. In dense complete settings, every ordered pair carries evidence, so forward-weight maximization aligns naturally with tournament and ranking formulations \cite{ACN08,CFR10}. In sparse settings, many pairwise relations are absent, heterogeneous, or structurally constrained by applications such as dependency repair or precedence resolution.

This is exactly why a method that performs well on sparse graph-benchmarks may still lose on LOLIB. The issue is not inconsistency in the experiments; it is a change in problem structure. Related work therefore has to prepare the reader for a comparative message that is strong but bounded.

This paper occupies a specific position in the literature. It is not a new approximation-theory paper in the style of local-ratio foundations \cite{BYGR98} or directed-feedback algorithm design \cite{DF03,ENSS98}. It is not an exact-method paper in the style of \cite{BSNA21}. It is not a web-scale systems paper in the style of \cite{SST16}. It is also not a dense linear-ordering or trainable ranking paper in the style of the LOLIB benchmark tradition \cite{MRD12}, tournament local search \cite{FLRS10}, or graph-neural ranking work such as GNNRank \cite{GNNRank22}. Instead, it is an integrated computational heuristic study whose main object is a reproducible framework for nonnegative weighted directed graphs.

That framework combines three roles that are often separated in the literature: a cycle-reduction seed grounded in prior local-ratio ideas, a complementary weighted seed, and incumbent-protected SCC neighborhood refinement. IPSNS is the main algorithmic contribution beyond the engineered LR-TA seed. It concentrates search on strongly connected components with positive backward contribution, applies SCC-local destroy-and-repair moves with restricted local-ratio repair, and accepts a candidate only when the global backward weight strictly decreases. The practical significance of that combination is twofold. First, it gives a transparent explanation for why the method should be robust on sparse directed graphs: the initial seeds construct feasible orders in different ways, and the refinement concentrates effort on strongly connected regions where ordering uncertainty remains. Second, it yields a claim that is strong enough to matter and narrow enough to defend: the framework is highly effective on the main sparse benchmark, near-optimal on the exact-validation subset, and explicitly limited on dense LOLIB instances where the external matrix-based pairwise-ordering baseline is stronger.
